\documentclass[11pt,a4paper]{article}
% This text is inserted in the beginning of all
% LaTex and Tex files I create.
%
% File created: Fri Jul 07 2023
% File name:    report_template.tex
% Path:         /home/name/Classes/AMPIII/Template/
%
% Tatiana Balbi Fraga
% Jul, 2023
%

\usepackage{babel}
\usepackage[utf8]{inputenc}

%recommended by fancyhdr package
\setlength{\headheight}{14.49998pt}
\addtolength{\topmargin}{-2.49998pt}

% include a minimal set of useful packages
\usepackage{graphicx}
\usepackage{amsfonts} 
\usepackage{amssymb}
\usepackage{amsmath}
\usepackage[a4paper,left=3.0cm,top=3.0cm,right=2.0cm,bottom=2.0cm]{geometry}
\usepackage{lastpage}
\usepackage{fancyhdr}
\usepackage{natbib,stfloats}
\usepackage[parfill]{parskip}
\usepackage{multirow}
\usepackage{helvet}

% PUT YOUR TITLE AND NAME HERE
\newcommand{\titlestr}{Planejamento do transporte coletivo em vans \\ Projeto PIBIC}
\newcommand{\shorttitlestr}{Planejamento do transporte coletivo em vans}
\newcommand{\authorstr}{Profa. Tatiana Balbi Fraga} % INSERT YOUR NAME(S)

\begin{document}
%%%%%%%%%%%%%%%%%%%%%%%%%%%%%%555
% title page
\begin{titlepage}
  \centering
  
  {\LARGE \titlestr \par}

  \vspace{1cm}
  {\Large \authorstr \par}

  {\bf Aluna: Dayane Eduarda da Silva}

  \vspace{1cm}
  \today     % PUT YOUR DATE HERE

  \vspace{1.5cm}
  Projeto submetido ao
  {\bf Edital PROPESQI no 02/2022}
  da Pró-Reitoria de Pesquisa e Inovação
  da Universidade Federal de Pernambuco

  \includegraphics[width=0.25\textwidth]{logos/ufpe_logo.png}

  \vspace{2cm}
  \flushleft
  {\bf Departamento de Engenharia de Produção} \\
  {\bf Núcleo de Tecnologia} \\
  {\bf Centro Acadêmico do Agreste} \\

  \vspace{5mm} {\footnotesize Este projeto é parte do projeto de pesquisa registrado na Propesqi intitulado 'Análise e Modelagem de Problemas de Otimização Contínua e Combinatória' \citep{Fraga2021}. Ao submeter este projeto eu declaro que o mesmo é de minha autoria e que todas as referências consultadas estão claramente citadas. O projeto está embasado no projeto desenvolvido pela aluna Dayane Eduarda da Silva, sob minha orientação, na disciplina 'Português Instrumental e Metodologia Científica' (disponível em ). Contudo o projeto que aqui apresento foi elaborado após extenso trabalho de pesquisa realizado por mim (incluindo levantamento bibliogŕafico e proposição da metodologia a ser aplicada).\par}
    
  \vspace{2mm} {\footnotesize O texto do projeto (atualizado), está disponível em forma de artigo no site MadScientistsDiary: https://github.com/tbfraga/gamos-madScientistsDiary-, portanto qualquer utilização, distribuição, replicação deste projeto deve estar em conformidade com a Licença Pública Internacional Creative Commons Atribuição-NãoComercial-SemDerivativos 4.0 e qualquer outro projeto desenvolvido com base neste projeto deve citá-lo como referência respeitando os direitos de autoria.\par}
    
  \vfill
\end{titlepage}

% put headings on each page
\pagestyle{fancy}
\fancyhf{}
\rhead{\shorttitlestr}
\lhead{\authorstr}
\rfoot{Page \thepage\ of \pageref{LastPage}}
\renewcommand{\headrulewidth}{1pt}


%%%%%%%%%%%%%%%%%%%%%%%%%%%%%%555
% main report
\clearpage
\section{Introdução}



\section{Fundamentação Teórica}



\section{Metodologia}

Este projeto será desenvolvido através de 5 atividades, conforme descrito a seguir:\\

Atividade 1: descrição detalhada sobre o processo de transporte coletivo estudado, informando toda a problemática envolvida (\emph{i.e.}, modo de planejamento do transporte atual, as incovêniencias, etc.).\\

Atividade 2: identificação de parâmetros para avaliação da qualidade do transporte.\\

Atividade 3: compreensão do modelo matemático e conceitual do problema a ser desenvolvido pela orientadora.\\

Atividade 4: solução do modelo e usando a ferramenta LINGO.\\

Atividade 5: levantamento de dados e comparação dos resultados antes e após aplicação da metodologia de solução desenvolvida.\\

Atividade 6: preparação de relatórios e artigo para o CONIC.

Tais atividades serão desenvolvidas pela aluna, sob a orientação e supervisão da Profa. orientadora (autora deste projeto).

\section{Resulados esperados}

Após a aplicação da metodologia acima descrita, esperamos esperamos obter os seguintes resultados: 

\begin{itemize}
    \item dados de históricos de vendas e outros dados relevantes referentes aos principais produtos da empresa estudada;
    \item mapeamento dos padrões de demanda destes produtos;
    \item aprofundamento do conhecimento dos participantes dos projetos sobre os comportamentos de vendas e sobre a identificação de padrões de demanda;
    \item artigo apresentado no CONIC.
\end{itemize}

Os resultados destes trabalhos também serão incluídos em pelo menos um artigo que será submetido para revista.

\section{Viabilidade de execução}

Os projetos serão realizados preferencialmente no CAA-UFPE. Assim, todos os participantes deste projeto, terão acesso garantido a toda a infraestrutura necessária para o correto desenvolvimento de seu trabalho incluindo recursos físicos (sala e mobiliário), bibliográficos e computacionais da UFPE e mais especificamente do CAA. O departamento de engenharia de produção do CAA conta atualmente com dois laboratórios de informática que disponibilizam, pelo menos, 30 computadores. O GAMOS, em especial, conta com laboratório próprio, e que atualmente dispõe de 3 computadores. Em termos de recursos bibliográficos, os pesquisadores da área de engenharia da produção contam com a biblioteca central da UFPE e as bibliotecas setoriais do CTG (Centro de Tecnologia e Geociências) e do CCEN (Centro de Ciências Exatas e da Natureza), localizadas no campus da UFPE de Recife, e com a biblioteca do próprio CAA, que possuem assinatura de alguns dos principais periódicos na área além do acesso remoto à base de dados disponíveis hoje via rede entre elas o banco de dados disponibilizados pela CAPES e pelo sciencedirect. 

Parte do projeto será também realizada através de visitas à empresa escolhida para realização desse trabalho. A parceria com a empresa será firmada durante a realização do projeto e o termo de parceria será anexado no relatório final.

\section{Cronograma de atividades do estudante}

Este projeto foi planejado para ser realizado durante o período de 1 ano, com início previsto para 2022/2023. As atividades descritas na metodologia estão projetadas para serem realizadas conforme cronograma apresentado na Tab. 3 :

\begin{table}[h]
\begin{center}
\begin{tabular}[c]{||c||c|c|c|c|c|c||}
\cline {1-7}
\multirow{2}{*}{Atividade} & \multicolumn{6}{c ||}{Cronograma (bimestre)} \\ \cline {2-7}
 & $1^o$ & $2^o$ & $3^o$ & $4^o$ & $5^o$ & $6^o$ \\ \cline {1-7}
$1^a$ & xx &  &  &  &  &  \\
$2^a$ & xx &  &  &  &  &  \\
$3^a$ &  & xx & xx &  &  &  \\
$4^a$ &  &  & xx & xx & xx &   \\
$5^a$ &  &  & xx & xx & xx & xx  \\
$6^a$ &  &  &  &  &  & xx  \\ \cline {1-7}
\end{tabular}
\label{tab:Cronograma}
\caption{Cronograma planejado para o projeto.}
\end{center}
\end{table}

%\bibliographystyle{plain}
%\bibliography{bib_file}

\clearpage
\begin{thebibliography}{12}

\bibitem[\protect\citeauthoryear{Ballou}{2001}]{Ballou2001}
Ballou, R.H. (2001).{\it Gerenciamento da Cadeia de Suprimentos: Planejamento, Organização e Logística Empresarial}, 4. ed., Porto Alegre: Bookman.

\bibitem[\protect\citeauthoryear{Ballou}{2006}]{Ballou2006}
Ballou, R.H. (2006).{\it Gerenciamento da Cadeia de Suprimentos / Logística Empresarial}, 5. ed., Porto Alegre: Bookman.

\bibitem[\protect\citeauthoryear{Boylan et al.}{2008}]{BoylanEtAl2008}
Boylan, J.E., Syntetos, A.A., e Karakostas, G.C. (2008). 'Classification for forecasting and stock control:a case study'. {\it Journal of the Operational Research Society}, Vol. 59, pp. 473--481.

\bibitem[\protect\citeauthoryear{Businger e Read}{1999}]{BusingerRead1999}
Businger, M.P., e Read, R.R. (1999). 'Identification of demand patterns for selective processing: acase study'. {\it Omega, Int. J. Mgmt Sci.}, Vol. 27, pp. 189--200.

\bibitem[\protect\citeauthoryear{Eaves e Kingsman}{2004}]{EavesKingsman2004}
Eaves A.H.C., e Kingsman B.G. (2004). 'Forecasting for the ordering and stock-holding of spare parts'. {\it J. O. Opl. Res. Soc.}, Vol. 55, pp. 431--437.

\bibitem[\protect\citeauthoryear{Fraga}{2021}]{Fraga2021}
Fraga, T.B. (2021). 'Análise e Modelagem de Problemas de Otimização Contínua e Combinatória'. Projeto de Pesquisa registrado em 21/09/2018, e aprovado pela Pró-reitoria de Pesquisa da UFPE em 09/12/2021 (Processo SIPAC 23076.037097/2021-67).

\bibitem[\protect\citeauthoryear{Makridakis et al.}{1998}]{MakridakisEtAl1998}
Makridakis, S.G.,Wheelwright, S.C., Hyndman, R.J. (1998).{\it Forecasting: Methods and Applications}, 3. ed., Wiley.

\bibitem[\protect\citeauthoryear{Makridakis e Hibon}{2000}]{MakridakisHibon2000}
Makridakis, S. e Hibon, M. (2000) 'The M3-Competition: results, conclusions and implications'. {\it International Journal of Forecasting}, Vol. 16, pp. 451--476.

\bibitem[\protect\citeauthoryear{Syntetos et al.}{2005}]{SyntetosEtAl2005}
Syntetos, A.A., Boylan, J.E., e Croston, J.D. (2005) 'On the categorization of demand patterns'. {\it Journal of the Operational Research Society}, Vol. 56 (5), pp. 495--503.

\bibitem[\protect\citeauthoryear{Ulrich et al.}{2022}]{UlrichEtAl2022}
Ulrich, M., Jahnke, H., Langrock, R., Pesch, R., e Senge, R. (2022) 'Classification-based model selection in retail demand forecasting'. {\it International Journal of Forecasting}, Vol. 38 (1), pp. 209--223.

\bibitem[\protect\citeauthoryear{Werner et al.}{2006}]{WernerEtAl2006}
Werner, L, Lemos, F.O., Daudt, T. (2006) 'Previsão de demanda e níveis de estoque uma abordagem conjunta aplicada no setor siderúrgico'. {\it XIII SIMPEP}, Bauru, SP, Brasil.

\bibitem[\protect\citeauthoryear{Williams}{1984}]{Williams1984}
Williams, T.M. (1984). 'Stock control with sporadic and slow-moving demand'. {\it Journal of the Operational Research Society}, Vol. 35 (10), pp. 939–948. 

\end{thebibliography} 

\end{document}
